\documentclass[conference]{IEEEtran}
\IEEEoverridecommandlockouts

% Base text: the author's draft "Cross-Language Static Analysis for ROS 2: Detecting Shadow
% Configuration Drift" (research_paper_robotics (1).pdf), reproduced unchanged.
% Every addition describing the components built since is marked with "% ADDED:".

\usepackage{cite}
\usepackage{amsmath,amssymb}
\usepackage{graphicx}
\usepackage{booktabs}
\usepackage{xcolor}
\usepackage{listings}
\usepackage{algorithm}
\usepackage{algpseudocode}
\usepackage{url}
\usepackage[hidelinks]{hyperref}

\lstset{basicstyle=\ttfamily\footnotesize,columns=fullflexible,keepspaces=true,
        frame=single,framerule=0.3pt,aboveskip=4pt,belowskip=4pt}
\newcommand{\code}[1]{\texttt{#1}}

\begin{document}

\title{Cross-Language Static Analysis for ROS 2:\\Detecting Shadow Configuration Drift}

\author{\IEEEauthorblockN{Soham Kishor Walam}
\IEEEauthorblockA{Robotics and Autonomous Systems Laboratory\\
Advanced Engineering Research\\
Email: soham.walam2005@gmail.com}}

\maketitle

\begin{abstract}
Modern ROS 2 robotics systems rely on complex multi-language configurations spanning XML
launch files, YAML parameters, Python scripts, and C++ implementations. Existing analysis tools
examine these files in isolation, missing critical inconsistencies that emerge through cross-file
interactions. We present rosfix, a cross-language static analysis tool that detects
\emph{Shadow Configuration Drift}---configuration inconsistencies visible only through
multi-file dependency analysis. Built using Tree-sitter AST parsing and unified intermediate
representation, rosfix analyzes 38,917 files across 12 major ROS 2 repositories, detecting 2,405
configuration issues with 87.3\% precision. Our evaluation demonstrates that 13.0\% of detected
issues are invisible to single-file analysis, validating the necessity of cross-language
approaches for ROS 2 configuration verification.
% ADDED:
We further describe a resolved launch graph that follows the semantics of the ROS~2 launch
system and \code{rcl}, per-distribution knowledge packs for installed packages, and a repair
engine that accepts a fix only after re-analysis proves it removes the finding; in a mutation
benchmark on five public repositories, all 107 injected faults were detected at the mutated line.
\end{abstract}

\begin{IEEEkeywords}
Static analysis, robotics middleware, configuration management, cross-language analysis
\end{IEEEkeywords}

% ===========================================================================
\section{Introduction}

The Robot Operating System 2 (ROS 2) has become the standard middleware for safety-critical
robotics applications through its Data Distribution Service (DDS) foundation and comprehensive
security framework. However, ROS 2's multi-language configuration model creates complex
dependencies between XML launch files, YAML parameters, Python launch scripts, and C++ source
code.

Consider a common scenario: an architect specifies \code{RELIABLE} QoS in a YAML configuration
file to ensure message delivery for safety-critical commands, but a developer hardcodes
\code{.best\_effort()} in the C++ node implementation to reduce latency. This creates what we
term \emph{Shadow Configuration Drift}---an inconsistency between intended configuration (YAML)
and runtime behavior (C++) that remains invisible to single-file analysis tools.

Existing ROS 2 analysis approaches focus on individual file types: QoS
Guard~\cite{lee2025qosguard} validates DDS XML profiles, LiSA4ROS2~\cite{zanatta2024lisa4ros2}
analyzes C++ source for security policy generation, and HAROS~\cite{santos2021haros} provides
ROS 1 code quality metrics. None address cross-file configuration consistency.

This paper contributes: (1) formal definition of Shadow Configuration Drift, (2) rosfix---a
cross-language static analysis framework detecting configuration inconsistencies across ROS 2
file types, (3) empirical evaluation on 38,917 files from 12 major repositories, and (4)
demonstration that 13.0\% of configuration issues require cross-file analysis for detection.

% ADDED:
In addition, it contributes (5) a resolved launch graph with \code{rcl}-compatible parameter-file
matching and launch-to-source linking (Sec.~\ref{sec:graph}); (6) installed-environment
awareness and per-distribution knowledge packs (Sec.~\ref{sec:installed}); (7) a verified repair
engine with a closed-set repair assistant and a sandboxed build check (Sec.~\ref{sec:repair});
and (8) a reproducible mutation benchmark at real sites (Sec.~\ref{sec:mutation}). Empirical
studies motivate these additions: Canelas et al.\ found that 27 of 50 sub-categories of ROS
misconfiguration were not addressed by any of 18 surveyed techniques~\cite{canelas2024misconfig},
and Chen et al.\ studied 121 interaction bugs across ten ROS projects~\cite{chen2025ibugs}.

% ===========================================================================
\section{Problem Formulation}

\subsection{Shadow Configuration Drift Definition}

Let $C = \{C_{xml}, C_{yaml}, C_{py}, C_{cpp}\}$ represent configuration sources in a ROS 2
system. The effective runtime configuration $C_{eff}$ results from composition with precedence
rules:
\begin{equation}
C_{eff} = C_{cpp} \oplus C_{py} \oplus C_{yaml} \oplus C_{xml}
\end{equation}
where $\oplus$ represents override composition (later sources override earlier ones).

Shadow Configuration Drift occurs when source code silently overrides high-level configuration
intent: Given dependency relation $D \subseteq C \times C$, Shadow Configuration Drift is the set
of inconsistencies detectable only through transitive closure $D^*$:
\begin{equation}
\mathit{ShadowDrift} = \Big\{ d \in \mathit{ConfigDrift} : d \notin \bigcup_i \mathit{SingleFile}(C_i)
\wedge d \in \mathit{CrossFile}(D^*) \Big\}
\end{equation}

\subsection{Motivating Examples}

\textbf{QoS Override}: YAML specifies \code{reliability: RELIABLE} for a sensor topic, but C++
node uses \code{rclcpp::QoS(10).best\_effort()}, silently degrading communication reliability
under network stress.

\textbf{Security Bypass}: XML governance requires \code{ENCRYPT} protection, but Python launch
script binds to localhost \code{127.0.0.1}, bypassing network-level security policies.

\textbf{Lifecycle Violation}: Launch file configures managed lifecycle for motor controller, but
C++ implementation publishes commands immediately without proper state transitions.

% ADDED:
\textbf{Ignored QoS override}: A parameter file sets
\code{qos\_overrides./scan.subscription.reliability: best\_effort}, but the subscription was
created without \code{QosOverridingOptions}; the override is ignored and the QoS in the code wins.
Detecting this requires the parameter file, the launch graph that binds it to the node, and the
node's source.

% ADDED:
\textbf{Namespaced parameter keys}: \code{rcl} matches parameter-file keys against fully
qualified node names, so a section keyed \code{amcl} is silently ignored for an AMCL node launched
in namespace \code{robot1} (\code{/robot1/amcl}).

% ===========================================================================
\section{Approach}

\subsection{Architecture Overview}

rosfix implements a four-stage analysis pipeline:
1.~\textbf{Multi-Language Parsing}: Language-specific parsers for XML (roxmltree), YAML
(serde\_yaml), Python (rustpython-parser), and C++ (Tree-sitter)
2.~\textbf{Unified IR Generation}: Cross-language intermediate representation capturing ROS 2
semantic elements
3.~\textbf{Dependency Resolution}: Cross-file dependency graph construction with parameter
propagation tracking
4.~\textbf{Shadow Drift Detection}: Pattern-based inconsistency detection across configuration
sources

\subsection{Cross-Language Intermediate Representation}

The unified IR captures ROS 2 abstractions independent of source file format:
\begin{lstlisting}
struct CrossLanguageIR {
  nodes: Vec<NodeDefinition>,
  parameters: Vec<ParameterDefinition>,
  topics: Vec<TopicDefinition>,
  qos_profiles: Vec<QoSProfile>,
  dependencies: DependencyGraph,
  file_origins: HashMap<ElementId, FileOrigin>
}
\end{lstlisting}
File origin tracking enables precise error reporting with source location information across all
configuration sources.

\subsection{Shadow Drift Detection Patterns}

We implement seven drift detection patterns:

\textbf{QoS Override}: Detects hardcoded QoS settings in source code that contradict
configuration file specifications

\textbf{Security Policy Bypass}: Identifies source code that circumvents SROS2 security
configurations

\textbf{Parameter Conflicts}: Finds competing parameter values across configuration layers

\textbf{Lifecycle Violations}: Validates managed lifecycle requirements against implementation

\textbf{Topic Type Mismatches}: Ensures publisher-subscriber type compatibility

\textbf{Executor Safety}: Detects blocking operations in single-threaded callbacks

\textbf{Physical Safety}: Validates URDF constraints and navigation footprint specifications

% ---------------------------------------------------------------------------
% ADDED: subsections D-I describe the components built since the draft.
% ---------------------------------------------------------------------------

\subsection{Fact Extraction}
\label{sec:facts}

Each file is reduced to typed facts. Extractors cover \code{package.xml}; a small CMake
interpreter (variables, executables and libraries, \code{install} rules, ament\_cmake\_auto,
rclcpp component registration, generate\_parameter\_library specifications);
\code{setup.py}/\code{setup.cfg}; launch files in XML, YAML and Python; parameter YAML; C++
via Tree-sitter~\cite{treesitter}; Python via the RustPython AST; URDF with a xacro subset
(properties, arguments, arithmetic, macros with block parameters, includes, conditionals); SROS2
governance, permissions and policy documents; and RViz configurations, whose displays carry
their own QoS.

Python launch files are not executed: a static interpreter evaluates the common subset
(\code{Node}, \code{IncludeLaunchDescription}, \code{GroupAction}, \code{PushRosNamespace},
\code{SetParameter}, \code{DeclareLaunchArgument}, \code{LaunchConfiguration}, path
substitutions, helper functions, \code{RewrittenYaml}/\code{ParameterFile}) and marks the rest
unknown. The parameter-YAML layer keeps exact byte spans and types scalars like
\code{rcl\_yaml\_param\_parser}, so that \code{1} is an integer, \code{1.0} a double and
\code{on} a boolean, exactly as at runtime. From nodes, the extractors record node names,
declared parameters with types and defaults, endpoints with literal names and statically
evaluated QoS, callbacks, callback groups, executors and blocking calls.

\subsection{Resolved Launch Graph}
\label{sec:graph}

Every launch file is a root. The resolver follows includes into workspace and installed
packages, propagates launch arguments (values passed by the includer override defaults),
composes namespace pushes (relative pushes compose, absolute ones reset), and applies remappings
by comparing fully expanded names. Parameter-file sections are matched like \code{rcl}: a key
$k$ applies to node $n$ iff
\begin{equation}
\mathrm{match}\big(\mathrm{abs}(k), \mathrm{fqn}(n)\big),\quad
\mathrm{abs}(k) = \begin{cases} k & k \text{ starts with } / \\ /k & \text{otherwise}\end{cases}
\end{equation}
with \code{*} matching one name token and \code{**} any number of tokens; for
\code{RewrittenYaml(root\_key=$r$)} keys are prefixed with $r$. Each launched node is linked to
its implementation through the package's build rules (executable sources, linked in-package
libraries, the same node classes in headers) or, for composable nodes, through the plugin class.

\subsection{Precision-First Analysis}
\label{sec:precision}

Every fact records what could not be determined, and every rule checks these markers before
reporting. Parameter declarations have three levels: \emph{closed} (all known),
\emph{dotted-open} (names containing a dot are uncertain, e.g.\ plugin parameters
\code{plugin\_name.param}) and \emph{open} (a fully dynamic name). Declarations made through
another object's node handle are attributed to a node class when possible; literal declarations
in the node's workspace dependency closure are accepted; a node deriving from a non-standard base
class that cannot be found is treated as open. Launch graphs with unresolved includes or dynamic
code are marked partial, and checks needing a closed world are suppressed for them.

\subsection{Installed Environment and Knowledge Packs}
\label{sec:installed}

In a sourced ROS shell, rosfix reads the ament resource index~\cite{ament_index} of every prefix
in \code{AMENT\_PREFIX\_PATH}: installed packages, executables (\code{lib/<pkg>}), pluginlib
plugins (\code{<category>\_\_pluginlib\_\_plugin}), rclcpp components and dependencies. Launch and
parameter files of installed packages included by the workspace are loaded for resolution; no
finding is reported inside them. Because binary installs contain no source, a per-distribution
\emph{knowledge pack}---rosfix's own extractors run once over the source of popular packages---
provides the parameters and QoS of installed nodes (Table~\ref{tab:packs}). Packs are used only
when the distribution is known, since parameter names change between distributions.

\begin{table}[t]
\centering
\caption{Embedded knowledge packs (built from nine upstream repositories)}
\label{tab:packs}
\footnotesize
\begin{tabular}{@{}lrrrrr@{}}
\toprule
Distribution & Packages & Executables & Components & Parameters & Size \\
\midrule
Humble & 40 & 100 & 52 & 2{,}596 & 11\,KiB\\
Jazzy  & 44 &  96 & 38 & 3{,}850 & 14\,KiB\\
Kilted & 44 &  97 & 37 & 3{,}711 & 14\,KiB\\
\bottomrule
\end{tabular}
\end{table}

\subsection{Rule Engine}
\label{sec:rules}

Nine domain agents run sequentially over the facts and the graph and report 44 rules with
stable identifiers (Table~\ref{tab:rules}). Each finding carries an exact location, related
locations, an explanation and, where possible, a fix.

\begin{table}[t]
\centering
\caption{Agents and rules}
\label{tab:rules}
\footnotesize
\begin{tabular}{@{}p{1.1cm}p{6.7cm}@{}}
\toprule
Agent & Rules \\
\midrule
params    & unmatched node key, missing \code{ros\_\_parameters}, undeclared parameter, type mismatch, missing required parameter (P001--P005)\\
topics    & near-miss names, remap without effect, type conflict, absolute name in namespace (T001--T005)\\
executor  & blocking wait/spin in a callback that shares the response's group or a single-threaded executor; sleep in callback (E001--E002)\\
control   & unknown joint, missing interface, controller-type typo, inconsistent config, \code{<ros2\_control>} joint not in URDF, missing hardware plugin (C001--C006)\\
packaging & executable not built/installed, undeclared dependency, \code{setup.cfg} script dir, uninstalled \code{launch/}, missing file (L001--L005)\\
simtime   & mixed \code{use\_sim\_time}, simulator without sim time, no clock (S001--S003)\\
qos       & reliability, durability, deadline/liveliness incompatibility, late joiner, ignored \code{qos\_overrides} (Q001--Q005)\\
nav2      & plugin section missing, plugin typo, unknown frame, footprint, inflation radius, parameter typo (N001--N006)\\
security  & weak governance, broad/expired permissions, unknown enclave, policy--graph mismatch, absolute keystore path (X001--X006)\\
\bottomrule
\end{tabular}
\end{table}

\subsection{Verified Repair}
\label{sec:repair}

Agents attach a fix only when the edit is local and mechanical (e.g.\ \code{1}$\to$\code{1.0},
renaming a misspelled key to the closest declared name, adding a \code{<depend>}). Fixes are
applied in memory, the whole workspace is re-analysed, and a fix is kept only if its finding
disappears and no new error appears (Algorithm~\ref{alg:fix}); accepted fixes are written
atomically, and only with an explicit \code{--write}.

\begin{algorithm}[t]
\caption{Fix verification}
\label{alg:fix}
\footnotesize
\begin{algorithmic}[1]
\State $F \gets$ non-overlapping fixes of all findings $D$
\For{round $=1..3$}
  \State $D' \gets \mathrm{analyse}(\text{workspace with } F \text{ applied in memory})$
  \State $R \gets \{f\in F \mid \mathrm{finding}(f)\in D' \lor \mathrm{newErrors}(D', D, f)\}$
  \If{$R=\emptyset$} \Return $F$ \EndIf
  \State $F \gets F \setminus R$
\EndFor
\State \Return $F$
\end{algorithmic}
\end{algorithm}

\textbf{Closed-set repair assistant.} Findings without a mechanical fix (e.g.\ a key
\code{max\_lin\_vel} for a node declaring \code{max\_linear\_velocity}) carry a repair site: the
exact span and the closed set of values valid there (declared parameters, \code{<ros2\_control>}
joints, URDF frames, executables, enclaves). The assistant chooses one value or abstains; it never
generates text, unlike generative repair models~\cite{silva2023repairllama,wang2023rapgen}. A
lexical ranker (edit distance, word overlap, ROS abbreviations) and a small learned ranker---a
16$\to$16$\to$1 perceptron over 16 string features, about 290 weights compiled into the binary---
score the candidates; a choice needs a clear margin, is verified by Algorithm~\ref{alg:fix}, and
is written only with an explicit flag, since re-analysis cannot prove intent. A prototype using a
0.5B-parameter code language model on the CPU needed 4.5\,s to load and 17.6\,s per finding,
which motivated the small ranker.

\textbf{Sandboxed build check.} With \code{--verify-build}, the workspace is copied twice to a
temporary directory and built before and after the fixes with \code{colcon}, in a ROS container
without network access when Docker is available; fixes that break a working build are refused.

\textbf{Optional references.} With \code{--online}, findings receive links to related Robotics
Stack Exchange questions and GitHub issues. Queries are built only from the rule that fired, so
no file content, path or name leaves the machine; a separate helper process performs the requests
against an HTTPS allowlist with timeouts and size limits, and results are cached.

\subsection{Implementation and Distribution}

rosfix is written in Rust (about 15{,}000 lines), runs single-threaded by default, and reports
text, JSON and SARIF~2.1.0~\cite{sarif} with exact regions, rule metadata and fixes, inline
suppressions, baselines for gradual adoption, and a non-zero exit code at a configurable severity.
Packaging is provided for Python wheels (with a \code{ros2 rosfix} command), standalone archives,
container images, a GitHub Action and pre-commit hooks.

% ===========================================================================
\section{Experimental Evaluation}

\subsection{Dataset and Methodology}

We evaluate rosfix on 12 major open-source ROS 2 repositories totaling 38,917 files:
PX4-Autopilot (26,404), Autoware (8,074), MoveIt2 (1,567), Navigation2 (1,333), rmf\_ros2 (695),
ros2\_control (314), rmw\_fastrtps (224), demos (217), system\_tests (104), turtlebot3 (88), sros2
(73), rmw\_cyclonedx (24).

All experiments execute on Windows 10.0.26200 with AMD Ryzen architecture using Rust 1.97.1
release builds.

\subsection{Detection Accuracy}

rosfix identifies 2,405 configuration issues across 947 files. Manual validation of a stratified
random sample ($n=312$, 95\% confidence) yields:
\begin{itemize}
\item Precision: 87.3\% (CI: 83.1\%--91.5\%)
\item Critical Issues: 127 (security, physical safety)
\item Error Level: 891 (functional problems)
\item Warning Level: 1,387 (performance, consistency)
\end{itemize}

\subsection{Shadow Drift Analysis}

Cross-file dependency analysis reveals 312 shadow drift instances (13.0\%) invisible to
single-file tools:
\begin{itemize}
\item QoS Overrides: 189 cases (hardcoded reliability downgrades)
\item Security Bypasses: 47 cases (policy circumvention)
\item Parameter Conflicts: 76 cases (competing values across layers)
\end{itemize}
Traditional single-file analyzers detect 0 instances of these cross-file inconsistencies.

\subsection{Performance Analysis}

\begin{table}[t]
\centering
\caption{Scalability and Parallelization Results}
\label{tab:scal}
\footnotesize
\begin{tabular}{@{}lrr@{}}
\toprule
Repository Size & Files & ms/File\\
\midrule
Small (demos) & 217 & 9.55\\
Large (MoveIt2) & 1,567 & 21.94\\
Very Large (Autoware) & 8,074 & 20.06\\
Complete Dataset & 38,917 & 13.81\\
\bottomrule
\end{tabular}
\end{table}

Rayon parallelization achieves 3.19$\times$ speedup on 4 cores and 5.15$\times$ on 8 cores, with
80\% and 64\% efficiency respectively.

% ---------------------------------------------------------------------------
% ADDED: subsections E-H report measurements of the current implementation.
% ---------------------------------------------------------------------------

\subsection{Resolved-Graph Analysis of Five Repositories}
\label{sec:five}

We ran the current implementation (release build, one thread, cold file cache, no ROS
installation) on shallow clones of five repositories and reviewed every finding manually
(Table~\ref{tab:five}). The L002 findings are undeclared \code{package.xml} dependencies, e.g.\
\code{turtlebot3\_bringup} starts nodes from \code{rclcpp\_components}, \code{camera\_ros} and
\code{image\_view} without declaring them, which breaks installation from binaries. E002 findings
are sleeps in executor callbacks; S002 findings are Nav2 example nodes started next to a simulator
without \code{use\_sim\_time}.

\begin{table}[t]
\centering
\caption{Current implementation on five repositories (cold cache, one thread)}
\label{tab:five}
\footnotesize
\begin{tabular}{@{}lrrrl@{}}
\toprule
Repository & Files & Time & Memory & Findings\\
\midrule
ros2/demos     & 138 & 1.3\,s & 13\,MB & E002$\times$2, L002$\times$6\\
ros2/sros2     &  20 & 0.1\,s &  5\,MB & ---\\
turtlebot3     &  63 & 0.1\,s & 12\,MB & E002$\times$2, L002$\times$5\\
ros2\_control  & 124 & 0.7\,s & 20\,MB & ---\\
navigation2    & 732 & 5.1\,s & 68\,MB & L002$\times$31, S002$\times$10\\
\bottomrule
\end{tabular}
\end{table}

Analysing real code also exposed false-positive sources that unit tests had not, each fixed in
the extractors or the resolver: parameter files passed through
\code{ParameterFile(RewrittenYaml(\dots))}; nested dictionaries in Python launch parameters
(flattened to dotted names by ROS); project-specific QoS classes such as
\code{LatchedPublisherQoS}; launch\_testing files; callbacks handing work to a thread;
metapackage dependencies; the member form of \code{declare\_or\_get\_parameter}; parameters
declared by shared library code in dependency packages; and template calls misparsed without
type information.

\subsection{Mutation Benchmark at Real Sites}
\label{sec:mutation}

A benchmark injects one fault at a time, in memory, at a site discovered in the clean analysis
(e.g.\ a parameter that a node really declares and a launch file really binds); a mutation counts
as detected only if the expected rule fires at the mutated line. All 107 mutations were detected
(Table~\ref{tab:mut}). Three faults injected by hand into Nav2's default parameter file (a
misspelled \code{controller\_frequency}, an integer for a double, a misspelled
\code{ros\_\_parameters}) were also detected, each with a verified fix. The operators are
synthetic and were written alongside the tool.

\begin{table}[t]
\centering
\caption{Mutation benchmark (five repositories)}
\label{tab:mut}
\footnotesize
\begin{tabular}{@{}llrr@{}}
\toprule
Mutation & Rule & Injected & Detected\\
\midrule
parameter key typo                   & P003 & 10 & 10\\
double written as integer            & P004 &  7 &  7\\
\code{ros\_\_parameters} misspelled  & P002 & 33 & 33\\
executable name typo in launch file  & L001 & 21 & 21\\
dependency removed from package.xml  & L002 & 16 & 16\\
Nav2 plugin type typo                & N002 & 20 & 20\\
\midrule
total                                &      & 107 & 107\\
\bottomrule
\end{tabular}
\end{table}

\subsection{Installed Packages}

With the Jazzy knowledge pack, a workspace launching Nav2's installed controller server with
\code{controller\_frequency: 20} is reported as a type error that crashes the node at startup
(fix: \code{20.0}), and a misspelled \code{min\_x\_velocity\_treshold} as undeclared---without
Nav2's source on the machine.

\subsection{Coverage}

In navigation2, rosfix built 50 launch graphs with 1,161 node instances, 1,065 of them linked to
source. Parameter checks apply fully to 132, to top-level names for 877 plugin-hosting nodes, and
are disabled for 56 nodes whose declarations cannot be determined.

% ===========================================================================
\section{Case Study: TurtleBot3}

Detailed analysis of the TurtleBot3 educational platform demonstrates rosfix's practical
effectiveness:

\textbf{Physical Safety Issues}: 6 URDF continuous joints missing velocity and effort limits,
potentially causing hardware damage in simulation or real deployment

\textbf{Executor Deadlocks}: 16 instances of blocking \code{sleep()} calls inside ROS 2 callbacks
that freeze single-threaded executors

\textbf{Policy Violations}: Multiple launch files missing \code{respawn="true"} attributes for
critical motor control nodes

All detected issues represent legitimate configuration problems confirmed through manual
inspection and testing.

% ADDED:
With the current implementation, TurtleBot3 yields seven findings (Table~\ref{tab:five}): two
sleeps inside action-server callbacks of \code{turtlebot3\_example} and five undeclared
dependencies (\code{rclcpp\_components}, \code{camera\_ros} and \code{image\_view} in
\code{turtlebot3\_bringup}; \code{rviz2} in the cartographer and navigation2 packages).

% ===========================================================================
\section{Related Work}

QoS Guard~\cite{lee2025qosguard} provides rigorous DDS QoS profile analysis but operates
exclusively on XML files, missing source code overrides that rosfix detects.

LiSA4ROS2~\cite{zanatta2024lisa4ros2} uses abstract interpretation for SROS2 policy generation
from source code but focuses on policy synthesis rather than configuration validation.

HAROS~\cite{santos2021haros} pioneered static analysis for ROS 1 but requires adaptation for
ROS 2's DDS-based architecture and QoS model.

rosfix uniquely addresses cross-file configuration consistency through unified multi-language
analysis, detecting issues invisible to existing single-domain approaches.

% ADDED:
ROSDiscover~\cite{timperley2022rosdiscover} statically recovers and checks the run-time
architecture of ROS~1 systems; rosfix resolves ROS~2 launch graphs, including parameter-file
binding, installed packages and packaging. Empirical studies of ROS
misconfigurations~\cite{canelas2024misconfig} and interaction bugs~\cite{chen2025ibugs} motivate
our rule selection. In program repair, small models are typically fine-tuned to generate
patches~\cite{silva2023repairllama} and benefit from retrieved examples~\cite{wang2023rapgen};
our assistant instead ranks a closed set of valid values and verifies every choice.

% ===========================================================================
\section{Limitations}

\textbf{Static Analysis Bounds}: Cannot detect runtime-only configuration issues or dynamic
parameter modifications

\textbf{Dependency Resolution}: Current implementation handles direct dependencies; complex
transitive relationships may be missed

\textbf{Ground Truth Scale}: Manual validation covers representative sample ($n=312$) rather than
all 2,405 detected issues

\textbf{Platform Dependencies}: Evaluation limited to Windows environment with specific hardware
configuration

% ADDED:
\textbf{Interpreter Subsets}: The Python-launch and xacro interpreters cover common subsets;
other parts are marked partial and dependent checks are suppressed.

% ADDED:
\textbf{Synthetic Mutations}: The mutation operators were written alongside the tool; the
benchmark shows that checks work on real code layouts, not recall on faults in the wild.

% ADDED:
\textbf{Knowledge Packs and Assistant}: Packs reflect the branches they were built from. The
learned ranker is trained on synthetic corruptions of real parameter names from the packs; on
7,351 held-out cases from packages not seen in training it chose the right value 3,816 times
versus 2,213 for the lexical ranker, with 132 versus 142 wrong edits. On injected mistakes in
\code{demos} and \code{navigation2} it fixed all 24 misspelled executable names (lexical: 11), but
both rankers abstained on all 20 misspelled parameter keys, although the right key was among the
options in 19 of them. These cases are synthetic rather than fixes from real histories, and the
sandboxed build check was not exercised in the evaluation (it requires Docker or colcon).

% ===========================================================================
\section{Conclusion}

This work introduces Shadow Configuration Drift as a formal characterization of cross-file
configuration inconsistencies in ROS 2 systems. Our rosfix framework demonstrates that
multi-language static analysis can detect 13.0\% more configuration issues than single-file
approaches, with 87.3\% precision across 38,917 files from major robotics projects.

The key insight is that ROS 2's multi-language configuration model requires corresponding
multi-language analysis approaches. Traditional tools examining XML, YAML, Python, or C++ files in
isolation miss critical inconsistencies that emerge through cross-file interactions.

Future work includes dynamic analysis integration to detect runtime-only issues, machine learning
enhancement for complex drift pattern recognition, and extended evaluation across additional
robotics domains.

% ADDED:
Beyond detection, the resolved launch graph, knowledge packs and verified repair move rosfix
toward use in continuous integration on arbitrary ROS~2 workspaces; next steps are a benchmark of
real historical fixes mined from permissively licensed repositories and training the learned
ranker on it.

\textbf{Availability}: rosfix is open source at
\url{https://github.com/SohamWalam11/ros2-safety-lint} with complete experimental datasets and
reproducible benchmarks.

\bibliographystyle{IEEEtran}
\bibliography{references}

\end{document}
